\documentclass[11pt,aspectratio=169]{ctexbeamer}

\usetheme{default}
\setbeamertemplate{navigation symbols}{}

\usepackage{amsmath}
\usepackage{booktabs}

\definecolor{DeepBlue}{RGB}{25,74,122}
\definecolor{SoftInk}{RGB}{53,59,66}
\definecolor{PaperBg}{RGB}{249,248,245}
\definecolor{PanelBlue}{RGB}{236,243,249}

\setbeamercolor{normal text}{fg=SoftInk,bg=PaperBg}
\setbeamercolor{structure}{fg=DeepBlue}
\setbeamercolor{frametitle}{fg=DeepBlue,bg=PaperBg}
\setbeamercolor{title}{fg=DeepBlue}
\setbeamercolor{block title}{fg=DeepBlue,bg=PanelBlue}
\setbeamercolor{block body}{fg=SoftInk,bg=white}

\newcommand{\Core}[1]{
  \begin{block}{核心结论}
    #1
  \end{block}
}

\newcommand{\LectureMeta}[2]{
  \title{大模型推理系统与实践}
  \subtitle{#1}
  \author{张书豪}
  \institute{华中科技大学计算机学院}
  \date{2027 · #2}
}

\newcommand{\PaperRef}[3]{
  \item \textbf{#1}（#2）\hfill\href{#3}{原文链接}
}

\newcommand{\ResearchSources}[1]{
  \begin{frame}{本讲参考资料}
    \small
    \begin{itemize}
      #1
    \end{itemize}
    \vspace{0.3em}
    \textcolor{SoftInk}{课件中的论文数字均指作者在特定模型、硬件和工作负载下报告的结果，不代表无条件性能。}
  \end{frame}
}

\newcommand{\InternalCaseNote}{
  \small\textcolor{SoftInk}{说明：本页为课程研究案例或教学假设，不代表已发表论文结论；只有在实验包给出模型、硬件、工作负载和原始记录时，数字才可作为结果引用。}
}

% 14 讲细分版的统一参考资料入口。论文数字的解释边界见课程内容基线。

\newcommand{\LectureOneSources}{\ResearchSources{
  \PaperRef{vLLM: Easy, Fast, and Cheap LLM Serving with PagedAttention}{SOSP 2023}{https://doi.org/10.1145/3600006.3613165}
  \PaperRef{SGLang: Efficient Execution of Structured Language Model Programs}{arXiv:2312.07104}{https://arxiv.org/abs/2312.07104}
  \PaperRef{nano-vllm-hust 教学代码}{课程仓库}{https://github.com/intellistream/intro-to-llm-inference-engines}
}}

\newcommand{\LectureTwoSources}{\ResearchSources{
  \PaperRef{vLLM / PagedAttention}{SOSP 2023}{https://doi.org/10.1145/3600006.3613165}
  \PaperRef{Sarathi-Serve}{OSDI 2024}{https://www.usenix.org/conference/osdi24/presentation/agrawal}
  \PaperRef{DistServe}{OSDI 2024}{https://www.usenix.org/conference/osdi24/presentation/zhong-yinmin}
}}

\newcommand{\LectureThreeSources}{\ResearchSources{
  \PaperRef{vLLM / PagedAttention}{SOSP 2023}{https://doi.org/10.1145/3600006.3613165}
  \PaperRef{DistServe}{OSDI 2024}{https://www.usenix.org/conference/osdi24/presentation/zhong-yinmin}
  \PaperRef{Llumnix}{OSDI 2024}{https://www.usenix.org/conference/osdi24/presentation/sun-biao}
}}

\newcommand{\LectureFourSources}{\ResearchSources{
  \PaperRef{Orca}{OSDI 2022}{https://www.usenix.org/conference/osdi22/presentation/yu}
  \PaperRef{Sarathi-Serve}{OSDI 2024}{https://www.usenix.org/conference/osdi24/presentation/agrawal}
  \PaperRef{FastServe}{NSDI 2026}{https://www.usenix.org/conference/nsdi26/presentation/wu-bingyang}
}}

\newcommand{\LectureFiveSources}{\ResearchSources{
  \PaperRef{vLLM / PagedAttention}{SOSP 2023}{https://doi.org/10.1145/3600006.3613165}
  \PaperRef{vAttention}{arXiv:2405.04437}{https://arxiv.org/abs/2405.04437}
  \PaperRef{Rethinking KV Cache Compression}{MLSys 2025}{https://proceedings.mlsys.org/paper_files/paper/2025/hash/26289c647c6828e862e271ca3c490486-Abstract-Conference.html}
  \PaperRef{Mooncake}{FAST 2025}{https://www.usenix.org/conference/fast25/presentation/qin}
}}

\newcommand{\LectureSixSources}{\ResearchSources{
  \PaperRef{DistServe}{OSDI 2024}{https://www.usenix.org/conference/osdi24/presentation/zhong-yinmin}
  \PaperRef{RAGCache}{arXiv:2404.12457}{https://arxiv.org/abs/2404.12457}
  \PaperRef{Mooncake}{FAST 2025}{https://www.usenix.org/conference/fast25/presentation/qin}
  \PaperRef{NEO}{MLSys 2025}{https://proceedings.mlsys.org/paper_files/paper/2025/hash/66a026c0d17040889b50f0dfa650e5e0-Abstract-Conference.html}
}}

\newcommand{\LectureSevenSources}{\ResearchSources{
  \PaperRef{vLLM / PagedAttention}{SOSP 2023}{https://doi.org/10.1145/3600006.3613165}
  \PaperRef{SGLang}{arXiv:2312.07104}{https://arxiv.org/abs/2312.07104}
  \PaperRef{DistServe}{OSDI 2024}{https://www.usenix.org/conference/osdi24/presentation/zhong-yinmin}
  \PaperRef{Mooncake}{FAST 2025}{https://www.usenix.org/conference/fast25/presentation/qin}
}}

\newcommand{\LectureEightSources}{\ResearchSources{
  \PaperRef{FlashAttention}{NeurIPS 2022}{https://arxiv.org/abs/2205.14135}
  \PaperRef{FlexGen}{ICML 2023}{https://proceedings.mlr.press/v202/sheng23a.html}
  \PaperRef{PowerInfer}{SOSP 2024 / arXiv:2312.12456}{https://arxiv.org/abs/2312.12456}
  \PaperRef{FlashInfer}{MLSys 2025}{https://proceedings.mlsys.org/paper_files/paper/2025/hash/dbf02b21d77409a2db30e56866a8ab3a-Abstract-Conference.html}
}}

\newcommand{\LectureNineSources}{\ResearchSources{
  \PaperRef{DEEPSERVE}{USENIX ATC 2025}{https://www.usenix.org/conference/atc25/presentation/hu-junhao}
  \PaperRef{QServe}{MLSys 2025}{https://proceedings.mlsys.org/paper_files/paper/2025/hash/fbe2b2f74a2ece8070d8fb073717bda6-Abstract-Conference.html}
  \PaperRef{vLLM-HUST}{开源维护仓库}{https://github.com/vLLM-HUST/vllm-hust}
}}

\newcommand{\LectureTenSources}{\ResearchSources{
  \PaperRef{vLLM / PagedAttention}{SOSP 2023}{https://doi.org/10.1145/3600006.3613165}
  \PaperRef{Revisiting Disaggregated LLM Serving}{2025 preprint, arXiv:2601.08833}{https://arxiv.org/abs/2601.08833}
  \PaperRef{Rethinking KV Cache Compression}{MLSys 2025}{https://proceedings.mlsys.org/paper_files/paper/2025/hash/26289c647c6828e862e271ca3c490486-Abstract-Conference.html}
}}

\newcommand{\LectureElevenSources}{\ResearchSources{
  \PaperRef{MLPerf Inference}{MLCommons benchmark}{https://mlcommons.org/benchmarks/inference-datacenter/}
  \PaperRef{ACM Artifact Review and Badging}{ACM policy}{https://www.acm.org/publications/policies/artifact-review-and-badging-current}
  \PaperRef{Revisiting Disaggregated LLM Serving}{2025 preprint, arXiv:2601.08833}{https://arxiv.org/abs/2601.08833}
}}

\newcommand{\LectureTwelveSources}{\ResearchSources{
  \PaperRef{nano-vllm-hust}{课程教学代码}{https://github.com/intellistream/intro-to-llm-inference-engines}
  \PaperRef{vLLM-HUST}{HUST 维护的 vLLM fork}{https://github.com/vLLM-HUST/vllm-hust}
  \PaperRef{上游 vLLM}{开源项目}{https://github.com/vllm-project/vllm}
}}

\newcommand{\LectureThirteenSources}{\ResearchSources{
  \PaperRef{课程内容基线与实验材料}{课程仓库}{https://github.com/intellistream/intro-to-llm-inference-engines}
  \PaperRef{DEEPSERVE}{USENIX ATC 2025}{https://www.usenix.org/conference/atc25/presentation/hu-junhao}
  \PaperRef{GhostServe}{MLSys 2026}{https://proceedings.mlsys.org/paper_files/paper/2026/hash/d99e8e80a6c41e148db686918dd7eab3-Abstract-Conference.html}
}}

\newcommand{\LectureFourteenSources}{\ResearchSources{
  \PaperRef{课程完整材料与后续更新}{课程仓库}{https://github.com/intellistream/intro-to-llm-inference-engines}
  \PaperRef{SpecGen}{MLSys 2026}{https://proceedings.mlsys.org/paper_files/paper/2026/hash/66a026c0d17040889b50f0dfa650e5e0-Abstract-Conference.html}
  \PaperRef{GhostServe}{MLSys 2026}{https://proceedings.mlsys.org/paper_files/paper/2026/hash/d99e8e80a6c41e148db686918dd7eab3-Abstract-Conference.html}
}}



\LectureMeta{第 13 讲 研究选题与课程项目工作坊}{研究生课程课件}

\begin{document}

\begin{frame}[plain]
  \titlepage
\end{frame}

\begin{frame}{教学目标}
  \Core{本讲围绕课程项目的问题定义、最小实现计划与验证路径展开。}
  \begin{itemize}
    \item 识别瓶颈对象和可控变量。
    \item 明确影响链路和验证方法。
    \item 为中期汇报和最终交付做收束。
  \end{itemize}
\end{frame}

\begin{frame}{项目定义的四个基本问题}
  \Core{若项目无法清楚回答以下四个问题，通常说明题目尚未形成稳定研究问题。}
  \begin{enumerate}
    \item 你要优化的系统对象是什么？
    \item 你真正改变的控制点是什么？
    \item 收益会沿着哪条链路传递到端到端指标？
    \item 什么结果会直接否定当前假设？
  \end{enumerate}
\end{frame}

\begin{frame}{常见项目方向}
  \Core{课程项目优先围绕真实 runtime、真实 benchmark 和可复现系统问题展开。}
  \begin{itemize}
    \item 调度：准入、尾时延控制、批处理组织。
    \item KV：块管理、复用、分层驻留。
    \item 状态：记忆对象、知识缓存与统一抽象。
    \item 执行：图执行、异构放置与数据路径优化。
    \item 工具：benchmark、可观测性与回归验证。
  \end{itemize}
\end{frame}

\begin{frame}{如何缩小问题范围}
  \Core{研究生项目中的常见问题在于题目定义过大，从而难以形成可信交付。}
  \begin{itemize}
    \item 先找一个最小可运行场景。
    \item 先做一个最小可验证机制。
    \item 先找到一个能够证伪自己的实验。
  \end{itemize}
\end{frame}

\begin{frame}{课程项目中的常见失控原因}
  \Core{项目中途失控通常源于问题对象、验证路径或最小交付边界未能及时收紧。}
  \begin{itemize}
    \item 问题太大，无法在一学期内完成。
    \item 假设太模糊，无法设计对照实验。
    \item 实现太重，挤占了分析和验证时间。
  \end{itemize}
\end{frame}

\begin{frame}{中期汇报建议结构}
  \Core{中期汇报最重要的是暴露问题定义和验证计划，而不是展示已经做了多少代码。}
  \begin{itemize}
    \item 问题背景与瓶颈对象
    \item 当前假设与拟议机制
    \item 实验设计与主要风险
    \item 需要老师和同学帮助判断的部分
  \end{itemize}
\end{frame}

\begin{frame}{从 tutorial 作业到课程项目}
  \Core{课程项目不应从零起步，前期 tutorial 作业应被视为项目准备阶段的能力训练。}
  \begin{itemize}
    \item Tutorial 1 到 5 用于建立 workload 指标判断、请求路径、调度、KV 与状态管理能力。
    \item Tutorial 6 到 10 用于建立架构、优化、平台、论文比较与实验方法能力。
    \item Tutorial 11 到 13 用于把开源系统实践、项目工作坊和最终汇报连接起来。
  \end{itemize}
\end{frame}

\begin{frame}{助教与学生的交付接口}
  \Core{如果课程采用 \texttt{\detokenize{nano-vllm-hust}} 统一作业包，学生提交与助教汇总可以保持同一结构。}
  \begin{itemize}
    \item 学生根据各个 \texttt{\detokenize{experiment_*/assignment_spring-2026/}} 中的模板提交实现、报告与结果。
    \item 助教可使用 \texttt{\detokenize{collect_benchmark_results.py}} 汇总 benchmark 结果。
    \item 统一目录结构有助于后续过渡到课程项目与最终汇报。
  \end{itemize}
\end{frame}

\begin{frame}{最终项目交付应该包含什么}
  \Core{最终项目不是一堆脚本，而应是一份可被他人理解和继续推进的系统工作样本。}
  \begin{itemize}
    \item 问题定义与机制设计说明。
    \item 最小实现或实验脚本。
    \item 结果分析与失败边界。
    \item 后续可扩展方向。
  \end{itemize}
\end{frame}

\begin{frame}{讲解：项目工作坊首先要把题目压缩成什么}
  \Core{一个课程项目如果不能在一句话里同时说出对象、控制点和指标，通常说明题目还没成形。}
  \begin{itemize}
    \item “对象”回答你到底在优化什么，例如调度、KV、状态搬运、观测路径。
    \item “控制点”回答你准备改哪里，而不是泛泛地说“优化系统”。
    \item “指标”回答你如何判断成功，以及什么结果会证明你判断错了。
  \end{itemize}
\end{frame}

\begin{frame}{讲解：为什么项目工作坊不鼓励一开始就做大而全}
  \Core{课程项目最常见的失败，并不是学生不努力，而是题目边界从第一周起就没有收紧。}
  \begin{itemize}
    \item 没有最小场景时，实验无法快速给出正反反馈。
    \item 没有最小机制时，实现工作会吞掉分析和验证时间。
    \item 所以工作坊的目标不是激励大家“做更大”，而是帮助大家“做得更真、更可证伪”。
  \end{itemize}
\end{frame}

\begin{frame}{例子：一个好题目和一个坏题目的差别（1）}
  \Core{项目工作坊最先要解决的不是“做什么酷”，而是“什么问题有清晰对象、控制点和验证路径”。}
\begin{block}{坏题目}
    “优化大模型推理性能”过大，既没有对象边界，也没有明确失败条件。
  \end{block}
\end{frame}

\begin{frame}{例子：一个好题目和一个坏题目的差别（2）}
  \Core{项目工作坊最先要解决的不是“做什么酷”，而是“什么问题有清晰对象、控制点和验证路径”。}
\begin{block}{好题目}
    “在长短请求混部场景下，验证某种 admission 策略能否改善短请求 TTFT 且不损伤 goodput”。
  \end{block}
\end{frame}

\begin{frame}{例子：一个好题目和一个坏题目的差别（3）}
  \Core{项目工作坊最先要解决的不是“做什么酷”，而是“什么问题有清晰对象、控制点和验证路径”。}
\begin{block}{课程结论}
    题目一旦能被写成“对象 + 控制点 + 指标 + 失败条件”，项目才真正开始成立。
  \end{block}
\end{frame}

\begin{frame}{例子：为什么中期汇报最该暴露风险而不是展示进度（1）}
  \Core{中期汇报的价值在于及时暴露错误方向，而不是把已经投入的工程量包装成确定性成果。}
\begin{block}{好汇报}
    说明当前假设最危险的地方、最可能失败的实验和最需要外界判断的路径。
  \end{block}
\end{frame}

\begin{frame}{例子：为什么中期汇报最该暴露风险而不是展示进度（2）}
  \Core{中期汇报的价值在于及时暴露错误方向，而不是把已经投入的工程量包装成确定性成果。}
\begin{block}{坏汇报}
    只列做了多少代码、跑了多少图，却回避核心假设是否站得住。
  \end{block}
\end{frame}

\begin{frame}{例子：为什么中期汇报最该暴露风险而不是展示进度（3）}
  \Core{中期汇报的价值在于及时暴露错误方向，而不是把已经投入的工程量包装成确定性成果。}
\begin{block}{教学意义}
    中期不是“半成品展示会”，而是“研究方向纠偏会”。
  \end{block}
\end{frame}

\begin{frame}{代码例子：项目提案的一页式模板（1）}
  \Core{工作坊里最好强制学生先写一页提案，再决定是否进入较重实现。}
\begin{block}{模板片段}
    {\ttfamily\small
    \detokenize{Problem object: ...}\par
    \detokenize{Changed control point: ...}\par
    \detokenize{Expected impact path: ...}\par
    \detokenize{Falsifying result: ...}\par
    }
  \end{block}
\end{frame}

\begin{frame}{代码例子：项目提案的一页式模板（2）}
  \Core{工作坊里最好强制学生先写一页提案，再决定是否进入较重实现。}
\begin{block}{课堂应追问}
    如果其中一项写不出来，说明项目还缺什么？
  \end{block}
\end{frame}

\begin{frame}{代码例子：项目提案的一页式模板（3）}
  \Core{工作坊里最好强制学生先写一页提案，再决定是否进入较重实现。}
\begin{block}{课程结论}
    一页提案的作用不是形式化，而是逼迫项目从“想法”进入“可检验命题”。
  \end{block}
\end{frame}

\begin{frame}{代码例子：最小实验矩阵如何设计（1）}
  \Core{课程项目不需要一开始就大规模扫参，更需要一个能快速否定错误方向的实验矩阵。}
\begin{block}{实验矩阵草图}
    {\ttfamily\small
    \detokenize{baseline vs method}\par
    \detokenize{short-heavy / mixed / long-heavy workload}\par
    \detokenize{ttft / p99 / goodput / probe metrics}\par
    }
  \end{block}
\end{frame}

\begin{frame}{代码例子：最小实验矩阵如何设计（2）}
  \Core{课程项目不需要一开始就大规模扫参，更需要一个能快速否定错误方向的实验矩阵。}
\begin{block}{课堂应追问}
    为什么先做三种 workload 的小矩阵，往往比直接追求“大而全”更有效？
  \end{block}
\end{frame}

\begin{frame}{代码例子：最小实验矩阵如何设计（3）}
  \Core{课程项目不需要一开始就大规模扫参，更需要一个能快速否定错误方向的实验矩阵。}
\begin{block}{结论}
    研究初期最宝贵的不是更多图，而是更快得到可否定反馈。
  \end{block}
\end{frame}

\begin{frame}{例子：项目失控往往从哪一步开始（1）}
  \Core{大多数课程项目不是在最后失败，而是在最初没有把边界收紧时就已经埋下失败。}
\begin{block}{失控点一}
    先承诺一个“大一统系统”，没有最小场景和最小机制。
  \end{block}
\end{frame}

\begin{frame}{例子：项目失控往往从哪一步开始（2）}
  \Core{大多数课程项目不是在最后失败，而是在最初没有把边界收紧时就已经埋下失败。}
\begin{block}{失控点二}
    先写很多代码，后补问题定义和实验设计。
  \end{block}
\end{frame}

\begin{frame}{例子：项目失控往往从哪一步开始（3）}
  \Core{大多数课程项目不是在最后失败，而是在最初没有把边界收紧时就已经埋下失败。}
\begin{block}{失控点三}
    只追求正结果，没有提前设计失败也有价值的解释路径。
  \end{block}
\end{frame}

\begin{frame}{课程项目题目池：KV 与状态方向}
  \Core{真实研究题目可以拆成适合课程项目的最小版本，重点是对象、控制点和验证路径清楚。}
  \begin{itemize}
    \item KV-Delta：比较 full-image handoff 与 delta handoff，解释 bytes 与 E2E 的关系。
    \item Segment Reuse：构造 raw negative 与 aligned positive，检查 commit ledger 与 correctness。
    \item Tiered KV：实现 movement log guard 或 restore reject reason summary。
  \end{itemize}
\end{frame}

\begin{frame}{课程项目题目池：执行路径与异构方向}
  \Core{执行路径项目不要求学生发明大系统，要求他们能用 profiler 和 trace 证明瓶颈在哪里。}
  \begin{itemize}
    \item MLP Materialization：用 profiler 判断单算子优化是否值得继续。
    \item Decode Selector：设计一页 claim ledger，区分 selection policy、routing regime 与端到端指标。
    \item KV-plane：整理机制、baseline 和 workload provenance，判断当前证据能否支撑论文主线。
  \end{itemize}
\end{frame}

\begin{frame}{课程项目题目池：调度与编排方向}
  \Core{调度与编排项目最容易夸大，因此必须提前写出 strongest baseline 和失败条件。}
  \begin{itemize}
    \item Topology-aware placement：比较 locality、pressure、session affinity 和 remote path 的取舍。
    \item Semantic MapReduce：区分 orchestration、coverage 与 reducer quality 三类 claim。
    \item shared workload provenance：为一个现有 benchmark 补 source、commit、adapter 和 failed-run 记录。
  \end{itemize}
\end{frame}

\begin{frame}{课程项目交付：一页 claim ledger}
  \Core{最终汇报不只交代码和图，还要交一页能让别人复查的 claim ledger。}
  \begin{itemize}
    \item 每个 claim 对应 artifact、baseline、workload、指标和设备。
    \item 每个 failed / blocked / invalid run 都要说明原因。
    \item 最后一行写：如果继续做，最小且最有价值的下一步是什么。
  \end{itemize}
\end{frame}

\begin{frame}{课堂练习：把你的题目压缩到一句话（1）}
  \Core{一个能做完的课程项目，往往都能被压缩成一句短而硬的问题句。}
\begin{block}{练习一}
    用一句话写出你的项目对象、控制点和主指标，不超过 35 个字。
  \end{block}
\end{frame}

\begin{frame}{课堂练习：把你的题目压缩到一句话（2）}
  \Core{一个能做完的课程项目，往往都能被压缩成一句短而硬的问题句。}
\begin{block}{练习二}
    再用一句话写出最危险的失败条件。
  \end{block}
\end{frame}

\begin{frame}{课堂练习：把你的题目压缩到一句话（3）}
  \Core{一个能做完的课程项目，往往都能被压缩成一句短而硬的问题句。}
\begin{block}{练习三}
    用一句话回答：如果只能保留一个实验，你会做哪个，为什么？
  \end{block}
\end{frame}


\begin{frame}{课堂讨论}
  \Core{课程项目首先应具备可证伪性，其次才适合继续优化。}
  \begin{itemize}
    \item 什么样的项目题目看起来大，但其实不具备验证路径？
    \item 为什么最小可运行场景对课程项目尤其重要？
    \item 中期汇报最应该暴露哪些不确定性？
  \end{itemize}
\end{frame}

\begin{frame}{本讲总结}
  \Core{工作坊的目标是将题目压缩为可落地、可验证的系统问题。}
  \begin{itemize}
    \item 项目成功的关键是问题定义和实验设计。
    \item 最后一讲将回到课程主线并完成最终汇报。
  \end{itemize}
\end{frame}

\begin{frame}{对应 Tutorial}
  \Core{本讲对应 Tutorial 12。先独立作答，再对照下一页参考答案。}
  \begin{itemize}
    \item 文件：\texttt{\detokenize{build/tutorials/Tutorial_12_课程项目工作坊.pdf}}
    \item 动手运行、日志记录与提交要求统一查看 \texttt{\detokenize{build/experiments/}} 和对应 experiment 源目录下的 \texttt{\detokenize{assignment_spring-2026/}}。
  \end{itemize}
\end{frame}

\end{document}